\subsection{二平面の小解消と射影的Kulikovモデル}
\label{subsec:grassmann-kulikov-v28}

前項の十八個の普通二重点は、二つの平面を使って同時に解消できる。
ここでいうKulikovモデルは、滑らかな全空間、被約単純正規交差特殊ファイバー、
自明な相対標準束を持つK3曲面の一径数退化を意味する
\cite{KulikovDegenerations,PerssonPinkhamDegenerations}。
以下では基底の $0$ を含む開近傍に制限し、他の特異ファイバーがあれば除いておく。

\begin{theorem}[平面を用いた射影的小解消]
\label{thm:grassmann-kulikov-v28}
系\ref{cor:eighteen-smoothing-nodes-v28}の一般一次摂動族を
$\mathscr X\to S$ とする。
特殊ファイバーの二成分 $\Pi_+,\Pi_-$ を $\mathscr X$ のWeil因子として扱い、
そのイデアルを順次blow-upすると、射影的な小解消
\[
 \widetilde{\mathscr X}\longrightarrow\mathscr X
\]
を得る。これはcrepantであり、$\widetilde{\mathscr X}$ は滑らかで
$\omega_{\widetilde{\mathscr X}/S}\simeq\mathcal O$ である。
特殊ファイバーは
\[
 \widetilde X=\widetilde\Pi_+\cup R\cup\widetilde\Pi_-,\qquad
 \widetilde\Pi_\pm\simeq
 \operatorname{Bl}_{Z_\pm}\mathbb P^2
\]
という三成分の被約単純正規交差曲面である。
ここで $Z_\pm$ は前項の平滑化切断の九個の相異なる零点であり、
$R$ および二重曲線 $D_\pm$ は同型を保つ。
この族は、両端が有理曲面、中間が楕円線織面であるType~IIのKulikovモデルとなる。
基底変換は必要ない。
\end{theorem}

\begin{proof}
普通二重点の近傍では、特殊ファイバーの二成分を $a=0$ と $b=0$ として、
全空間の方程式は
\[
 ab-tH(a,b,c,t)=0,\qquad H_c(0)\ne0
\]
と書ける。$c$ を $H$ で置き換えれば、基底パラメータ $t$ と二成分を保ったまま
\begin{equation}\label{eq:node-small-resolution-v28}
 ab=tc
\end{equation}
という局所形になる。
平面側を $\{a=t=0\}$ とし、そのイデアル $(a,t)$ をblow-upする。
$a$ が生成元となるチャートでは
\[
 t=au,\quad b=uc
\]
であり、$(a,u,c)$ が滑らかな三次元座標を与える。
特殊ファイバーは $au=0$ である。
$t$ が生成元となるチャートでは
\[
 a=tv,\quad c=vb
\]
となり、$(t,v,b)$ が滑らかな座標を与える。
この二チャートを合わせると、元の普通二重点だけが一本の射影直線に置き換わる。
従ってこれは小解消であり、解消後の特殊ファイバーは被約かつ単純正規交差である。

選択した成分 $\{a=t=0\}$ への制限は、その局所座標 $(b,c)$ の原点のblow-upである。
実際、上の二チャートでは $b=uc$ と $c=vb$ がその二つの標準チャートとなる。
他方の成分 $\{b=t=0\}$ は $(a,c)$ を座標とする $u=0$ のチャートで同型に保たれる。
従って平面を選んで解消すれば、平面だけが九点でblow-upされ、$R$ は変わらない。

全空間の普通二重点以外では、各 $\Pi_\pm$ はCartier因子である。
平面の内部では特殊ファイバー自身を局所的に定義する因子であり、
二重曲線の平滑化切断が非零の点では $ab=t\cdot\text{unit}$ より一つの式で定義される。
従ってこのblow-upは普通二重点の外で同型である。
二平面は交わらないため、それぞれのイデアルのblow-upは互いの中心に影響しない。
これにより、局所小解消は大域的かつ射影的な二回のblow-upとして実現する。

$\mathscr X$ はGrassmann多様体と基底の積の中で、
$\bigl(\bigwedge^2\mathcal U^*\bigr)^{\oplus3}$ の正則切断により定義される。
相対随伴により $\omega_{\mathscr X/S}\simeq\mathcal O$ である。
小解消には例外因子がなく、$\mathscr X$ の標準束は可逆なので、
標準因子の差は零となる。従って解消はcrepantであり、
$\omega_{\widetilde{\mathscr X}/S}\simeq\mathcal O$ を得る。
一般ファイバーは変更されず、種数二十のK3曲面のままである。
上のチャートが示す $au=t$ 型の局所形により、解消後の族は半安定である。
特殊ファイバーは有理曲面、楕円線織面、有理曲面の鎖となり、Type~IIの形を持つ。
\end{proof}

\begin{corollary}[二重曲線の法束の相殺]
\label{cor:normal-cancellation-v28}
解消後の二重曲線について
\[
 N_{D_\pm/\widetilde\Pi_\pm}
 \simeq N_{D_\pm/R}^{-1},\qquad
 \mathcal T^1_{\widetilde X}|_{D_\pm}\simeq\mathcal O_{D_\pm}
\]
が成立する。両端の有理曲面では $D_\pm$ は自己交差数零の反標準因子である。
\end{corollary}

\begin{proof}
摂動の一次平滑化切断を $h_\pm$ とする。
その零除数は $Z_\pm$ であり、前項から
\[
 \mathcal O_{D_\pm}(Z_\pm)
 \simeq N_{D_\pm/R}\otimes\mathcal O_{D_\pm}(3)
\]
である。滑らかな三次曲線上の九点をblow-upすると、その法束は
\[
 N_{D_\pm/\widetilde\Pi_\pm}
 =\mathcal O_{D_\pm}(3)(-Z_\pm)
 =N_{D_\pm/R}^{-1}
\]
となる。これが局所一次変形層の自明性も与える。
また $\operatorname{Bl}_9\mathbb P^2$ 上では、三次曲線の狭義変換は
$3H-\sum_{j=1}^9E_j=-K$ であり、その自己交差数は $9-9=0$ である。
\end{proof}

\paragraph{偏極と有理端成分の意味.}
Pl\"ucker偏極の引戻しは、解消後も半豊富で、一般ファイバー上で次数三十八の偏極を与える。
例外的な十八本の射影直線上では次数零となる。
また、端成分が九点blow-upで反標準曲線の自己交差数が零であることだけから、
反標準線形系が楕円pencilをなすとはいえない。
ここでは二つの端を有理曲面と記述し、法束の自明性や捩れ性を追加で仮定しない。
